
\documentclass{article}
\usepackage{colortbl}
\usepackage{makecell}
\usepackage{multirow}
\usepackage{supertabular}

\begin{document}

\newcounter{utterance}

\centering \large Interaction Transcript for game `eqbench', experiment `eq\_bench', episode 57 with reflexion{-}Qwen3.5{-}4B.
\vspace{24pt}

{ \footnotesize  \setcounter{utterance}{1}
\setlength{\tabcolsep}{0pt}
\begin{supertabular}{c@{$\;$}|p{.15\linewidth}@{}p{.15\linewidth}p{.15\linewidth}p{.15\linewidth}p{.15\linewidth}p{.15\linewidth}}
    \# & \multicolumn{2}{c}{Player} && \multicolumn{2}{c}{Game Master} \\
    \hline

    \theutterance \stepcounter{utterance}  
    & & & \multicolumn{4}{p{0.6\linewidth}}{
        \cellcolor[rgb]{0.9,0.9,0.9}{
            \makecell[{{p{\linewidth}}}]{
                \texttt{\tiny{[P1$\langle$GM]}}
                \\ 
\texttt{Your task is to predict the likely emotional responses of a character in this dialogue:} \\
\\ 
\texttt{Claudia: Tyler, you look well. Teaching suits you.} \\
\texttt{Tyler: Yeah, well, someone's gotta do the dirty work.} \\
\texttt{Claudia: I didn't mean it like that. I...} \\
\texttt{Tyler: I know what you meant. You always do, don't you?} \\
\texttt{Claudia: Not always, no. But I try to understand, at least.} \\
\texttt{Tyler: Understand? Like you understood when you left me alone?} \\
\texttt{Claudia: That's not fair, Ty. I had to...} \\
\texttt{Tyler: What? Chase your dreams? Leave your past behind?} \\
\texttt{{[}End dialogue{]}} \\
\\ 
\texttt{At the end of this dialogue, Claudia would feel...} \\
\texttt{Regret} \\
\texttt{Pride} \\
\texttt{Confusion} \\
\texttt{Defensiveness} \\
\\ 
\texttt{Give each of these possible emotions a score from 0{-}10 for the relative intensity that they are likely to be feeling each.} \\
\\ 
\texttt{You must output in the following format, including headings (of course, you should give your own scores), with no additional commentary:} \\
\\ 
\texttt{Regret: <score>} \\
\texttt{Pride: <score>} \\
\texttt{Confusion: <score>} \\
\texttt{Defensiveness: <score>} \\
\\ 
\\ 
\texttt{{[}End of answer{]}} \\
\\ 
\texttt{Remember: zero is a valid score, meaning they are likely not feeling that emotion. You must score at least one emotion > 0.} \\
\\ 
\texttt{Your answer:} \\
            }
        }
    }
     \\ \\

    \theutterance \stepcounter{utterance}  
    & \multicolumn{4}{p{0.6\linewidth}}{
        \cellcolor[rgb]{0.561,0.737,0.561}{
            \makecell[{{p{\linewidth}}}]{
                \texttt{\tiny{[P1$\rangle$GM]}}
                \texttt{Regret: 8} \\
\texttt{Pride: 2} \\
\texttt{Confusion: 1} \\
\texttt{Defensiveness: 9} \\
\\ 
\texttt{Claudia's response "That's not fair, Ty. I had to..." is a classic defensive reaction to Tyler's accusation that she left him to chase her dreams. She is defending her actions and trying to justify her departure, which indicates high defensiveness. While she might feel some regret for the pain she caused, her immediate reaction in the dialogue is to push back against Tyler's interpretation rather than express sorrow. She does not appear proud of leaving, nor is she confused about what happened; she clearly remembers her reasons.} \\
            }
        }
    }
    & & \\ \\

    \theutterance \stepcounter{utterance}  
    & & & \multicolumn{2}{p{0.3\linewidth}}{
        \cellcolor[rgb]{0.9,0.9,0.9}{
            \makecell[{{p{\linewidth}}}]{
                \texttt{\tiny{[GM$|$GM]}}
                \texttt{Regret: 8} \\
\texttt{Pride: 2} \\
\texttt{Confusion: 1} \\
\texttt{Defensiveness: 9} \\
            }
        }
    }
    & & \\ \\

    \theutterance \stepcounter{utterance}  
    & & & \multicolumn{2}{p{0.3\linewidth}}{
        \cellcolor[rgb]{0.9,0.9,0.9}{
            \makecell[{{p{\linewidth}}}]{
                \texttt{\tiny{[GM$|$GM]}}
                \texttt{Regret: 8} \\
\texttt{Pride: 0} \\
\texttt{Confusion: 0} \\
\texttt{Defensiveness: 6} \\
            }
        }
    }
    & & \\ \\

    \theutterance \stepcounter{utterance}  
    & & & \multicolumn{2}{p{0.3\linewidth}}{
        \cellcolor[rgb]{0.9,0.9,0.9}{
            \makecell[{{p{\linewidth}}}]{
                \texttt{\tiny{[GM$|$GM]}}
                \texttt{game\_result = LOSE} \\
            }
        }
    }
    & & \\ \\

    \theutterance \stepcounter{utterance}  
    & & \multicolumn{4}{p{0.6\linewidth}}{
        \cellcolor[rgb]{0.95,0.95,0.95}{
            \makecell[{{p{\linewidth}}}]{
                \texttt{\tiny{[P1$\leftrightarrow$P]}}
                \texttt{Post{-}episode Reflexion (not sent to GM):} \\
\texttt{New plan: I misinterpreted the emotional nuance of Claudia's final line. While her immediate reaction to Tyler's accusation was defensive ("That's not fair"), the context of the entire dialogue—where she admits she "had to" leave and Tyler explicitly calls out her selfishness—suggests that the underlying, lingering emotion driving her behavior and her admission of necessity is regret. I incorrectly prioritized the surface{-}level defensiveness over the deeper, more complex emotional state implied by her justification. Next time, I will analyze the full trajectory of the conversation, noting that her admission of having "to" leave (implying a difficult choice) combined with her attempt to explain rather than deny, points more strongly to Regret as the dominant feeling, potentially lowering the score for Defensiveness or balancing it against Regret based on the depth of her explanation. I will ensure I distinguish between a temporary defensive reaction and the core emotional state the character is experiencing.} \\
            }
        }
    }
    & \\ \\

\end{supertabular}
}

\end{document}
